\subsection{Diskrete Zufallsvariablen}
\shortdefinition Z.V. $\cX$ ist \bi{diskret}, falls endl. oder abzählb. Menge $W \subsetneq \R$ existiert, s.d. $\P[\cX \in W] = 1$ (Werte v. $\cX$ f.s. in $W$)

\shortremark Falls $\Omega$ endl. / abzählb., dann ist jede Z.V. $\cX$ diskret.

\shortdefinition[Gewichtsf.] $\cX$ mit $W$ endl.: $p_{\cX}(x) = \P[\cX = x]$

\shortdefinition[Verteilung] $(p(x))_{x \in W} \defEquiv \forall x \in W \; p(x)$

\shorttheorem $(p(x))_{x \in W} = \sum_{x \in W} p(x) = 1$

\shortremark $\forall (p(x))_{x \in W} \; \exists$ Z.V. mit dieser Verteilung. Können desh. schreiben:
``Sei $\cX$ disk. Z.V. mit Verteilung $(p(x))_{x \in W}$''


\subsubsection{Zusammenhang Verteilung, Verteilungsfunktion}
\shorttheorem $\cX$ disk. Z.V. wie oben, dann ist \hl{Verteilungsfunktion}: $\forall x \in \R \; F_\cX(x) = \sum_{\elementstack{y \in W}{y \leq x}} p(y)$.
\textbf{Umgekehrt:} $p(x)$ ist die ``Sprunghöhe'' im Punkt $x \in W$, $W$ pos. Sprünge in $F_\cX$
